\documentclass[11pt,letterpaper]{article}

% --- Packages ---
\usepackage[top=0.55in, bottom=0.55in, left=0.85in, right=0.85in]{geometry}
\usepackage{fontawesome5}
\usepackage{xcolor}
\usepackage[hidelinks]{hyperref}
\usepackage{enumitem}
\usepackage{parskip}
\usepackage{titlesec}
\usepackage{setspace}

% --- Colors ---
\definecolor{CursorPurple}{HTML}{5D3FD3}
\definecolor{DarkGrey}{HTML}{333333}

\setstretch{1.0}

\begin{document}

% --- Header ---
\begin{center}
    {\Huge \textbf{\color{CursorPurple} Aditya Kulkarni}}\\[8pt]
    {\color{DarkGrey} \small
    \faPhone \ (281) 876-7066 \quad \textbar \quad
    \faEnvelope \ \href{mailto:adityakulkarnius@gmail.com}{adityakulkarnius@gmail.com} \\[4pt]
    \faLinkedin \ \href{https://linkedin.com/in/adityakulkarni244}{linkedin.com/in/adityakulkarni244} \quad \textbar \quad
    \faGithub \ \href{https://github.com/Aditya-k24}{github.com/Aditya-k24}
    }
\end{center}

\vspace{-4pt}
{\color{CursorPurple}\rule{\linewidth}{1.5pt}}
\vspace{10pt}

% --- Date & Salutation ---
June 18, 2026\\[10pt]
\textbf{To the Cursor Team,}

\vspace{6pt}

% --- Opening ---
The highest-leverage thing an AI tool can do for an engineering team is remove the gap between what an engineer knows how to do and what they can actually ship. Cursor does that better than anything else I have used --- and I use it daily. The Forward Deployed Engineer role is about going inside customer teams to make that gap disappear at scale. That is the exact intersection of technical depth and customer impact I want to build my career at.

\vspace{4pt}

\begin{itemize}[leftmargin=*, itemsep=3pt, label={\color{CursorPurple}\faAngleRight}]
    \item \textbf{Evals, Tracing, and Model Behavior:} I built an LLM evaluation harness from scratch --- 2,000 frozen test examples, hallucination rate, p50/p95 latency, cost-per-call, McNemar significance tests across 6 metrics, comparing Claude, GPT-4o, and Gemini. When FDE work requires debugging why a customer's Cursor-powered workflow is degrading, knowing how to instrument and measure AI behavior is the difference between guessing and knowing.

    \item \textbf{Production AI Shipped to Enterprise Clients:} At \textbf{Isomer AI}, I shipped 10+ production LLM features --- multi-model pipelines, Temporal.io agentic workflows --- to enterprise customers, across 923 commits and 7 production deployments. I have written the code that customer engineers inherit and have to maintain. I understand what they need before they say it.

    \item \textbf{Agentic Systems, Prototype to Production:} \textbf{AlgoPulse} was an agentic AI orchestration platform I took from idea to 2,000 req/sec in production --- Temporal.io, Kafka, event-driven SSE delivery. The FDE motion is to ship production-grade systems within days. I have done exactly that on my own, and I want to do it embedded inside teams with real constraints.

    \item \textbf{Observability Infrastructure:} \textbf{CortexQ} has end-to-end observability I built --- OpenTelemetry, Prometheus, Grafana, per-stage latency tracking, circuit breaker tracing. Latency and cost tradeoff analysis is not theoretical for me; it is instrumented and measured in every system I build.
\end{itemize}

\vspace{6pt}

% --- Closing ---
Cursor is making the best engineers faster and the rest of the world capable of engineering for the first time. I want to be in the rooms where that happens. I would be glad to discuss.

\vspace{12pt}

\textbf{Sincerely,}\\[8pt]
Aditya Kulkarni

\end{document}
